До сих пор речь шла о случайных процессах со значениями в $\R$,
но при обсуждении марковских процессов и марковских цепей удобнее говорить о процессах с более общим пространством состояний $X$, наделенном $\sigma$-алгеброй $\mathcal{X}$. Тогда процесс состоит из измеримых отображений $\xi_t: (\Omega, \B) \to (X, \X)$. Хотя пространства состояний из реальных задач борелевски изоморфны прямой, они обычно обладают какой-то естественной структурой, которую невыгодно терять. Например, при
рассмотрении случайного блуждания по целочисленной решетке на
плоскости было бы неудобно отождествлять ее с $\Z$.

Пусть $T \subset \R$.
Как и ранее, для всякого процесса $\{\xi_t\}_{t \in T}$ появляются порожденные им $\sigma$-алгебры $\F_{=t} = \sigma(\xi_t)$, а также более широкие $\sigma$-алгебры
$\F_{\leq t} = \sigma(\xi_s : s \leq t)$ и $F_{\geq t} = \sigma(\xi_s : s \geq t)$, интерпретируемые как «прошлое» и «будущее».

\begin{definition}
    Процесс $\{\xi_t\}_{t \in T}$ называется марковским, если для всякого $t \in T$ всех $A \in \F_{\leq t}$ и $B \in \F_{\geq t}$ почти наверное 
    
    \[ P(A \cap B|\F_{=t}) = P(A|\F_{=t})P(B|\F_{=t}). (15.1) \]
    
    Как обычно, тождество для множеств равносильно тождеству
    \[ \E(fg|\F_{=t}) = \E(f|\F_{=t})\E(g|\F_{=t}) \]
    для ограниченных $\F_{\geq t}$-измеримых функций $f$ и ограниченных $\F_{\leq t}$-измеримых функций $g$.
\end{definition}

Смысл равенства — условная независимость будущего и прошлого при фиксированном настоящем. Для таких процессов будущее зависит от прошлого только через настоящее. Приведем еще одно равносильное условие. Как обычно, равенство условных средних есть равенство почти всюду.

\begin{lemma}
    Тождество (15.1) равносильно тождеству
    \[ P(B|\F_{\leq t}) = P(B|\F_{=t}) \quad \forall B \in \F_{\geq t}. (15.2) \]
    Иначе говоря,
    
    \[ \E(f|\F_{\leq t}) = \E(f|\F_{=t}) (15.3) \]
    
    для всякой ограниченной $F_{\geq t}$-измеримой функции $f$
\end{lemma}

\begin{proof}
    Если верно (15.1), то для всякого $A \in \F_{\leq t}$ интеграл от $I_AI_B$, всегда равный интегралу от $P(A \cap B|\F_{=t})$, есть интеграл от $P(A|F_{=t})P(B|\F_{=t})$, что равно интегралу от $I_A P(B|\F_{=t})$ из-за $\F_{=t}$-измеримости $P(B|\F_{=t})$. Это дает (15.2).
    
    Пусть выполнено (15.2). Надо проверить, что для $A \in \F_{\leq t}$ левая и правая части (15.1) дают равные интегралы по множествам $C \in \F_{=t}$. Интеграл левой части равен $P(A \cap B \cap C)$, ибо $I_C P(A \cap B|\F_{=t}) = P(A \cap B \cap C|F_{=t})$ в силу включения $C \in \F_{=t}$. Интеграл от правой части равен (также с учетом включения $C \in \F_{=t}$)

    \[ \int I_C P(A|\F_{=t})P(B|\F_{=t}) dP = \int P(A|\F_{=t})P(C \cap B|\F_{=t}) dP = \int I_A P(C \cap B|\F_{=t}) dP = \int I_A P(C \cap B|\F_{\leq t}) dP. \]

    Так как $A \in \F_{\leq t}$, то это тоже равно $P(A \cap B \cap C)$    
\end{proof}

В основном определении и равносильном описании из предыдущей леммы использовались множества из исходного вероятностного
пространства или функции на нем. В следующей характеризации марковости рассматриваются функции на фазовом пространстве $E$ (которое никак не фигурировало в определении) или множества из $E$.

\begin{theorem}
    Случайный процесс $\xi_t$ со значениями в $E$ является марковским в точности тогда, когда для всякой ограниченной $\X$-измеримой функции $g$ на $E$ при всех $s, t \in T, s \leq t$ верно равенство
    
    \[ \E(g(\xi_t)|\F_{\leq s}) = E(g(\xi_t)|\F_{=s}). (15.4) \]

    Равносильное условие $P(\xi_t \in E|\F_{\leq s}) = P(\xi_t \in E|\F_{=s}) \forall E \in \X$.
\end{theorem}

\begin{proof}
    Если верно (15.3), то верно и (15.4), ибо функция $f = g(\xi t)$ измерима относительно $F_{\geq s}$ при $t \geq s$.

    Пусть верно (15.4). Покажем, что верно и (15.3). Сначала рассмотрим функции $f$ вида $f = f_1(\xi_{t_1})\ldots f_n(\xi_{t_n}))$, где $t \leq t1 < t2 < \ldots < t_n$ и функции $f_i$ на $X$ измеримы относительно $\X$. Проверим (15.3) индукцией по $n$. При $n = 1$ это есть (15.4). Пусть равенство верно при некотором $n$, $t_n < t_{n + 1}$, функция $f_{n+1}$ ограничена и $\X$-измерима. Тогда 
    
    \[ \E((f_1(\xi_{t_1}) \ldots f_{n+1}(\xi_{t_{n + 1}} )|\F_{\leq t}) = \E(\E(f_1(\xi_{t_1}) \ldots f_{n + 1}(\xi_{t_{n + 1}})|\F_{\leq t_n})|\F_{\leq t}) = \]
    \[ = \E(f_1(\xi_{t_1})\ldots  f_n (\xi_{t_n}) \E (f_{n + 1}(\xi_{t_{n + 1}} )|\F_{\leq t_n})|\F_{\leq t}), \]
    
    ибо функции $f_j(\xi_{t_j})$ измерима относительно $\F_{\leq t_n}$ при $j \leq n$. В силу (15.4) мы имеем 
    
    \[ \E(f_{n + 1}(\xi_{t_{n + 1}} )|\F_{\leq t_n}) = \E(f_{n + 1}\xi_{t_{n + 1}} )|\F_{=t_n}), \] 
    
    причем функция справа имеет вид $\varphi(\xi_{t_n})$, как показано в лемме ниже.
    
    По предположению индукции
    \[ E(f_1(\xi_{t_1})\ldots f_n(\xi_{t_n})\varphi(\xi_{t_n})|\F_{\leq t}) = \E(f_1(\xi_{t_1})\ldots f_n(\xi_{t_n})\varphi(\xi_{t_n})|\F_{=t}). \]
    
    С этой же функцией совпадает $\E(f_1(\xi_{t_1})\ldots f_{n + 1}(\xi_{t_{n + 1}} )|F_{=t})$, так как $\F_{=t} \subset \F_{\leq t_n}$ и потому
    
    \[ \E(f_1(\xi_{t_1})\ldots f_{n + 1}(\xi_{t_{n + 1}} )|\F_{=t})= \E(\E(f_1(\xi_{t_1})\ldots f_{n + 1}(\xi_{t_{n + 1}} )|\F_{\leq t_n})|\F_{=t}) = \]
    
    \[ = \E(f_1(\xi_{t_1})\ldots f_n(\xi_{t_n})\E(f_{n + 1}(\xi_{t_{n + 1}} )|\F_{\leq t_n})|F_{=t}). \]
    
    Итак, равенство (15.4) верно для функций указанного вида, а тогда и для их линейных комбинаций.

    Наконец, для всякой ограниченной $F_{\geq t}$-измеримой функции $f$ найдется последовательность таких линейных комбинаций $f_j$, сходящаяся к $f$ в $L^2(P)$. Соответствующие условные средние тоже сходятся в $L^2(P)$, поэтому равенство (15.4) остается в силе и для $f$
\end{proof}

\begin{theorem}
    Пусть $S$~--- борелевское множество в $\R^n$ с борелевской $\sigma$-алгеброй $\mathcal{S}$ (или измеримое пространство, изоморфное такому множеству) и $\{\xi_t\}_{t \in T}$~--- марковский процесс с пространством состояний $S$, причем $T$~--- какое-либо из множеств $[0, +\infty)$, $(-\infty, +\infty)$, $\N$, $\Z$. Тогда этот процесс обладает переходной функцией, т.е. имеется функция $P(s, x, t, E)$ четырех аргументов, где $s, t \in T, x \in S$, $E \in \mathcal{S}$, для которой
    \begin{itemize}
        \item[(i)] при фиксированных $s, t, x$ функция $E \mapsto P(s, x, t, E)$ является вероятностной мерой на $\mathcal{S}$,
        \item [(ii)] при фиксированных $s, t, E$ функция $x \mapsto P(s, x, t, E)$ измерима относительно $\mathcal{S}$,
        \item[(iii)] $E \mapsto P(s, x, s, E) = \delta_x(E)$ есть мера Дирака в точке x,
        \item[(iv)] выполнено уравнение Чэпмена – Колмогорова

        \[ P(s, x, t, E) = \int_S P(u, y, t, E)P(s, x, u, dy), s \leq u \leq t, (15.5) \]
        \item[(v)] при фиксированных $s, t, E$ ($s \leq t$) почти наверное
        
        \[ P(s, \xi_s, t, E) = P(\xi_t \in E|\F_{=s}). \]
    \end{itemize}
\end{theorem}

\begin{note}
    Переходная функция дает техническую интерпретацию «вероятности процесса попасть в множество E в момент t при нахождении в точке $x$ в момент времени $s$».
    Из леммы вытекает равенство
    \[ P(s, \xi_s, t, E) = P(\xi_t \in E|\F_{\leq s}). \]
    Условие (v) можно записать в виде
    \[ P(s, x, t, E) = P(\xi_t \in E|\xi_s = x) \]
    В случае счетного множества состояний и $P(\xi_s = x) > 0$ правая часть понимается обычным образом, а в общем случае она понимается как $\theta(x)$, где $\theta$ — борелевская функция, которая представляет условную вероятность в виде $P(\xi_t \in E|\sigma(\xi_s)) = \varphi(\xi_s)$.
    Смысл равенства в (15.5) том, что попадание в $E$ в момент $t$ при нахождении в $x$ в момент $s$ слагается из переходов в различные $y$ в промежуточный момент $u$ и последующих переходов из $y$ в $E$.
\end{note}

\begin{definition}
    Если переходные функции $P(s, x, t, E)$ зависят только от разности $t - s$, то процесс называют однородным.
\end{definition}

\begin{theorem}
    Процесс с независимыми приращениями является марковским.
\end{theorem}

\begin{corollary}
    Винеровский и пуассоновский процессы являются марковскими
\end{corollary}